\documentclass[12pt, a4paper]{article}
\usepackage[utf8]{inputenc}
\usepackage[a4paper, margin=1.5cm]{geometry}
\usepackage{titlesec}
\usepackage{enumitem}
\usepackage{booktabs}
\usepackage{tabularx}
\usepackage{array}
\usepackage[colorlinks=true, linkcolor=blue, urlcolor=blue, citecolor=blue]{hyperref}
\usepackage{parskip}

% Clean section formatting
\titleformat{\section}{\large\bfseries}{\thesection.}{0.5em}{}
\titleformat{\subsection}{\normalsize\bfseries}{\thesubsection}{0.5em}{}

\begin{document}

% ==========================================
% COMP2300 OFFICIAL TITLE PAGE
% ==========================================
\begin{titlepage}
    \centering
    \vspace*{3cm}
    
    {\Huge \textbf{COMP2300: Software Design and Development Project} \par}
    \vspace{1.5cm}
    
    {\LARGE \textbf{Deliverable 1 (D1): Comprehensive Literature Review \& Technical Execution Report} \par}
    \vspace{0.5cm}
    {\Large \textit{Topic 1: Carbon Footprint Awareness and Reduction Tool (India-Centric Challenge)} \par}
    
    \vspace{2.5cm}
    
    \begin{tabular}{ll}
        \textbf{Document Version:} & Version 1.0 (Detailed Master Reference) \\
        \textbf{Date:} & October 18, 2026 \\
        \textbf{Module:} & COMP2300 \\
    \end{tabular}
    
    \vspace{2.5cm}
    
    {\large \textbf{Authors:} \par}
    \vspace{0.5cm}
    \begin{tabular}{ll}
        \textbf{Arnav} & \texttt{[arnav@soton.ac.uk]} \\
        \textbf{Hiranya} & \texttt{[hiranya@soton.ac.uk]} \\
        \textbf{Keshav} & \texttt{[keshav@soton.ac.uk]} \\
    \end{tabular}
    
    \vfill
    {\normalsize University of Southampton \par}
\end{titlepage}

% ==========================================
% SECTION 1: CORE CONCEPTS & METHODOLOGIES
% ==========================================
\section{Core Concepts and Research Methodologies (Beginner-Friendly Guide)}

To design an 80+ Distinction-level Carbon Footprint (CF) tool tailored to urban India, it is essential to establish clear definitions of the underlying environmental accounting and computer science methodologies used across the literature:

\begin{itemize}[leftmargin=*]
    \item \textbf{Direct vs. Indirect (Embodied) Emissions:} Direct emissions include physical fuel burned in vehicles (Scope 1) and coal-grid electricity consumed by home appliances like air conditioners and geysers (Scope 2). Indirect or ``embodied'' emissions (Scope 3) represent the hidden carbon released across the supply chain to produce goods and services---such as growing cotton for apparel or manufacturing single-use plastic containers for food delivery \cite{shi2021global, koide2019carbon}.
    \item \textbf{Life Cycle Assessment (LCA --- Bottom-Up Accounting):} A method that calculates the exact environmental impact of a specific product or activity from raw material extraction to disposal \cite{shi2021global, liu2020packaging}. In software, LCA is used when physical units are known (e.g., multiplying kilometers traveled by a vehicle's fuel emission factor).
    \item \textbf{Input-Output Analysis (IOA --- Top-Down Accounting):} A macroeconomic method that estimates carbon emissions based on financial expenditure or broad product categories when granular physical tracking is impractical \cite{shi2021global}.
    \item \textbf{Scientometric Review:} A quantitative meta-study that analyzes thousands of published research papers to identify macro shifts in scientific focus over time \cite{shi2021global}.
    \item \textbf{Spatial Micro-Data Analysis:} A methodology that links door-to-door household expenditure surveys with supply-chain databases to map how carbon footprints vary across specific districts and cities \cite{lee2021india}.
    \item \textbf{Quasi-Experimental Study \& Food Frequency Questionnaire (FFQ):} A study design that measures participant behavior before and after an educational or software intervention. Using an FFQ (scoring diet frequency on a simple 0--4 scale) avoids forcing users to weigh food in grams, preventing ``tracker fatigue'' \cite{malan2020impact}.
    \item \textbf{The Value-Action Gap:} A psychological phenomenon where individuals express high concern about climate change in surveys (values) but fail to alter daily high-emission habits (actions) due to a lack of immediate, quantitative feedback at the moment of decision-making \cite{khandelwal2019students}.
    \item \textbf{Cluster Analysis (K-Means) \& Gap Analysis:} An unsupervised machine learning method that automatically groups users with similar consumption patterns into distinct ``lifestyle profiles'' and calculates the exact gap between their current emissions and a sustainable target budget \cite{koide2019carbon}.
\end{itemize}

% ==========================================
% SECTION 2: DEEP-DIVE INTO THE 8 PAPERS
% ==========================================
\section{Detailed Analysis of the 8 Foundational Research Papers}

\subsection{1. Shi \& Yin (2021) --- Global Carbon Accounting \& Embodied Emissions}
\begin{itemize}[leftmargin=*]
    \item \textbf{Citation \& Link:} S. Shi and J. Yin (2021), \textit{Environmental Impact Assessment Review} \cite{shi2021global}. Link: \href{https://doi.org/10.1016/j.eiar.2021.106571}{https://doi.org/10.1016/j.eiar.2021.106571}
    \item \textbf{Core Learnings:} Establishes LCA (bottom-up) and IOA (top-down) as the two primary accounting pillars. Demonstrates that global CF research has transitioned away from heavy industry toward urban household consumption and the ``embodied carbon'' hidden in consumer services.
    \item \textbf{Research Gap Addressed:} Standard calculators only measure direct fuel and electricity, overlooking indirect retail and service emissions.
    \item \textbf{Software Execution Plan:} Justifies implementing a hybrid \texttt{EmissionCalculator} interface in Java using the \textbf{Strategy Pattern}. Direct activities (bus travel, AC usage) use bottom-up LCA math, while indirect habits (fast-fashion apparel, quick-commerce orders) use top-down IOA category multipliers.
\end{itemize}

\subsection{2. Lee, Taherzadeh, \& Kanemoto (2021) --- Indian Household Carbon Drivers}
\begin{itemize}[leftmargin=*]
    \item \textbf{Citation \& Link:} J. Lee et al. (2021), \textit{Global Environmental Change} \cite{lee2021india}. Link: \href{https://doi.org/10.1016/j.gloenvcha.2020.102205}{https://doi.org/10.1016/j.gloenvcha.2020.102205}
    \item \textbf{Core Learnings:} Analyzed micro-consumption data from 203,313 households across 623 Indian districts. Proved that high-expenditure urban households emit nearly 7 times more carbon than low-expenditure households, and that footprints vary heavily by district, local power grids, and regional dietary habits rather than Western car-ownership baselines.
    \item \textbf{Research Gap Addressed:} Proves empirically why Western calculators fail in India and why city-level localization is mandatory.
    \item \textbf{Software Execution Plan:} Powers our \textbf{3-Step Onboarding Wizard}. Selecting a base city (e.g., Dehradun, Gurugram, New Delhi) queries an immutable Java \texttt{HashMap} ($O(1)$ lookup) containing regional Indian grid emission factors and baseline budgets.
\end{itemize}

\subsection{3. Liu et al. (2020) --- Urban Food Delivery \& Packaging Waste}
\begin{itemize}[leftmargin=*]
    \item \textbf{Citation \& Link:} G. Liu et al. (2020), \textit{Waste Management} \cite{liu2020packaging}. Link: \href{https://doi.org/10.1016/j.wasman.2020.07.028}{https://doi.org/10.1016/j.wasman.2020.07.028}
    \item \textbf{Core Learnings:} Conducted a big-data LCA on urban food delivery services, proving that single-use polypropylene plastic containers, bags, and foils generate a massive Global Warming Potential (GWP) overhead alongside last-mile delivery exhaust.
    \item \textbf{Research Gap Addressed:} Validates our \textbf{Quick-Commerce \& Micro-Delivery} niche. Urban Indian users frequently order 10-minute groceries (Zepto, Blinkit) and gravy-heavy restaurant meals (Paneer Butter Masala, Dal Makhani) packaged in heavy plastic containers, which existing calculators ignore.
    \item \textbf{Software Execution Plan:} When a user logs a delivery, the \texttt{QuickCommerceCalculator} applies a specific packaging and two-wheeler last-mile penalty multiplier on top of the food's baseline footprint.
\end{itemize}

\subsection{4. Malan et al. (2020) --- Low-Friction Dietary Scoring \& Self-Efficacy}
\begin{itemize}[leftmargin=*]
    \item \textbf{Citation \& Link:} H. Malan et al. (2020), \textit{Nutrients} \cite{malan2020impact}. Link: \href{https://doi.org/10.3390/nu12092890}{https://doi.org/10.3390/nu12092890}
    \item \textbf{Core Learnings:} A quasi-experimental study ($n = 176$ students) showing that replacing tedious gram-by-gram food weighing with a simplified 0--4 composite score increased student self-efficacy and reduced dietary carbon footprints by 14\%.
    \item \textbf{Research Gap Addressed:} Exposes both a UX gap (manual weighing causes app abandonment) and a cultural gap (Western studies focus on beef, whereas Indian vegetarian/flexitarian diets are driven by high-emission dairy like paneer, butter, and ghee).
    \item \textbf{Software Execution Plan:} Replaces nested drop-down menus with a single smart-text input field backed by a lightweight \textbf{Naive Bayes classifier}, converting phrases like \textit{``Paneer lunch and abcoffee delivery''} into a composite Dairy \& Packaging score.
\end{itemize}

\subsection{5. Khandelwal et al. (2019) --- Indian Youth Psychology \& Shared Mobility}
\begin{itemize}[leftmargin=*]
    \item \textbf{Citation \& Link:} M. Khandelwal et al. (2019), \textit{Management and Economics Research Journal} \cite{khandelwal2019students}. Link: \href{https://doi.org/10.18639/MERJ.2019.735006}{https://doi.org/10.18639/MERJ.2019.735006}
    \item \textbf{Core Learnings:} Surveyed students at a leading Indian private university (Amity University). Found a strong ``value-action gap'': students care about climate change but lack quantitative awareness of daily emission drivers. Once informed of their personal impact, they show a high willingness to adopt shared transit.
    \item \textbf{Research Gap Addressed:} Validates targeting Indian university students/young professionals and highlights the lack of tools modeling regional \textbf{Tier-2 to Tier-1 intercity transit}.
    \item \textbf{Software Execution Plan:} Justifies our \textbf{Intercity Transit Intercept} module. When a user plans a trip (e.g., Dehradun to Gurugram), the JavaFX UI dynamically compares a solo cab against an AC diesel sleeper bus (FlixBus/Zingbus) and electrified rail (Vande Bharat + Metro), awarding badges for choosing shared rail/bus.
\end{itemize}

\subsection{6. Lim \& Hayder (2019) --- Campus Baselines \& HVAC Dominance}
\begin{itemize}[leftmargin=*]
    \item \textbf{Citation \& Link:} M. Lim and G. Hayder (2019), \textit{International Journal of Engineering and Advanced Technology} \cite{lim2019performance}. Link: \href{https://doi.org/10.35940/ijeat.A2672.109119}{https://doi.org/10.35940/ijeat.A2672.109119}
    \item \textbf{Core Learnings:} Evaluated campus operational footprints across electricity, transport, water, and waste, proving that electrical HVAC cooling/heating and daily commuting dominate student footprints.
    \item \textbf{Research Gap Addressed:} Justifies prioritizing high-load Indian seasonal appliances (ACs, desert coolers, bathroom geysers) paired with coal-heavy regional power grids.
    \item \textbf{Software Execution Plan:} Dedicated \textbf{Appliance Load} slider and \textbf{Hostel/Flat Benchmark} card on the dashboard, translating 1 hour of reduced AC/geyser usage into tangible monthly carbon budget savings.
\end{itemize}

\subsection{7. Brewer et al. (2015) --- HCI Gamification for Behavioral Shifting}
\begin{itemize}[leftmargin=*]
    \item \textbf{Citation \& Link:} R. S. Brewer et al. (2015), \textit{ACM CHI PLAY '15} \cite{brewer2015tough}. Link: \href{https://doi.org/10.1145/2793107.2793108}{https://doi.org/10.1145/2793107.2793108}
    \item \textbf{Core Learnings:} Demonstrated that static energy dashboards fail to sustain user attention, whereas casual game mechanics (points, streaks, visual feedback) actively help users understand and shift residential electricity habits (complemented by Kjeldskov et al.'s 2015 CHI work on predictive \textit{Eco-Forecasting} \cite{kjeldskov2015ecoforecasting}).
    \item \textbf{Research Gap Addressed:} Bridges the gap between static environmental spreadsheets and engaging Human-Computer Interaction (HCI).
    \item \textbf{Software Execution Plan:} Uses JavaFX reactive properties (\texttt{DoubleProperty}) styled with \textbf{AtlantaFX} to animate a color-shifting Live Carbon Gauge, weekly streaks, and time-series forecasting.
\end{itemize}

\subsection{8. Koide et al. (2019) --- Lifestyle Clustering \& Gap Analysis (ML Backbone)}
\begin{itemize}[leftmargin=*]
    \item \textbf{Citation \& Link:} R. Koide et al. (2019), \textit{Sustainability} \cite{koide2019carbon}. Link: \href{https://doi.org/10.3390/su11215983}{https://doi.org/10.3390/su11215983}
    \item \textbf{Core Learnings:} Applied cluster analysis and multiple regression to 47,000 households. Proved that over 50\% of emissions are indirect and that generic advice fails; clustering consumers into specific lifestyle segments enables targeted gap analysis.
    \item \textbf{Research Gap Addressed:} Statistical clustering has rarely been embedded into a live, personal desktop application for Indian households.
    \item \textbf{Software Execution Plan:} Embeds a lightweight \textbf{K-Means Clustering Engine} in Java. Weekly user logs are classified into profiles (e.g., \textit{``High Quick-Commerce Consumer''} vs. \textit{``Frequent Intercity Commuter''}) to dynamically serve hyper-personalized reduction nudges on the dashboard sidebar.
\end{itemize}

% ==========================================
% SECTION 3: SYNTHESIS TABLE
% ==========================================
\section{Synthesis Matrix: Literature to Software Architecture}

\begin{table}[h!]
\centering
\small
\renewcommand{\arraystretch}{1.3}
\begin{tabularx}{\textwidth}{>{\bfseries}p{3.2cm} p{3.8cm} X}
\toprule
Paper & Core Research Domain & Direct Feature in Our Java/JavaFX Tool \\
\midrule
Shi \& Yin \cite{shi2021global} & Macro Accounting (LCA/IOA) & Hybrid \texttt{EmissionCalculator} Strategy Pattern for direct \& indirect logs. \\
Lee et al. \cite{lee2021india} & Indian District Drivers & Onboarding City Selector mapped to regional grid \texttt{HashMaps}. \\
Liu et al. \cite{liu2020packaging} & Delivery Packaging Waste & Quick-Commerce Penalty Multiplier for Zepto/Blinkit \& gravy tubs. \\
Malan et al. \cite{malan2020impact} & Student Dietary Footprints & Low-friction NLP smart-text logging \& Indian Dairy composite score. \\
Khandelwal et al. \cite{khandelwal2019students} & Indian Youth Mobility & Intercity Transit Intercept (Solo Cab vs. Sleeper Bus vs. Electric Rail). \\
Lim \& Hayder \cite{lim2019performance} & Campus HVAC Baselines & Appliance Load sliders (AC, cooler, geyser) \& hostel cohort benchmarks. \\
Brewer et al. \cite{brewer2015tough} & HCI Gamification & Reactive AtlantaFX Carbon Gauge, streak counters, and milestone badges. \\
Koide et al. \cite{koide2019carbon} & Consumer Clustering & Live K-Means clustering engine serving profile-specific sidebar nudges. \\
\bottomrule
\end{tabularx}
\caption{Mapping the 8 foundational papers directly to COMP2300 software features.}
\end{table}

% ==========================================
% REFERENCES (ACM SIGCHI NUMBERED FORMAT)
% ==========================================
\begin{thebibliography}{10}

\bibitem{shi2021global}
Shaoqing Shi and Jianhua Yin. 2021. Global research on carbon footprint: A scientometric review. \emph{Environmental Impact Assessment Review} 89 (2021), 106571. DOI: \href{https://doi.org/10.1016/j.eiar.2021.106571}{https://doi.org/10.1016/j.eiar.2021.106571}

\bibitem{lee2021india}
Jemyung Lee, Oliver Taherzadeh, and Keiichiro Kanemoto. 2021. The scale and drivers of carbon footprints in households, cities and regions across India. \emph{Global Environmental Change} 66 (2021), 102205. DOI: \href{https://doi.org/10.1016/j.gloenvcha.2020.102205}{https://doi.org/10.1016/j.gloenvcha.2020.102205}

\bibitem{liu2020packaging}
Gengyuan Liu, Feni Agostinho, Huabo Duan, Biagio F. Giannetti, Hao Wang, Cecilia M. V. B. Almeida, Marco Casazza, and Zhifeng Yang. 2020. Environmental impacts characterization of packaging waste generated by urban food delivery services. A big-data analysis in Jing-Jin-Ji region (China). \emph{Waste Management} 117 (2020), 157--169. DOI: \href{https://doi.org/10.1016/j.wasman.2020.07.028}{https://doi.org/10.1016/j.wasman.2020.07.028}

\bibitem{malan2020impact}
Hannah Malan, G. A. Challamel, D. Silverstein, C. Hoffs, E. Spang, S. A. Pace, B. L. R. Malague\~{n}o, C. D. Gardner, M. C. Wang, W. Slusser, and J. A. Jay. 2020. Impact of a Scalable, Multi-Campus ``Foodprint'' Seminar on College Students' Dietary Intake and Dietary Carbon Footprint. \emph{Nutrients} 12, 9 (2020), 2890. DOI: \href{https://doi.org/10.3390/nu12092890}{https://doi.org/10.3390/nu12092890}

\bibitem{khandelwal2019students}
Manvi Khandelwal, V. Jain, A. Sharma, and S. Bansal. 2019. Students' Attitude toward Carbon Footprints of a Leading Private University in India. \emph{Management and Economics Research Journal} 5 (2019), 7606. DOI: \href{https://doi.org/10.18639/MERJ.2019.735006}{https://doi.org/10.18639/MERJ.2019.735006}

\bibitem{lim2019performance}
M. Lim and G. Hayder. 2019. Performance and reduction of carbon footprint for a sustainable campus. \emph{International Journal of Engineering and Advanced Technology} 9, 1 (2019), 3489--3493. DOI: \href{https://doi.org/10.35940/ijeat.A2672.109119}{https://doi.org/10.35940/ijeat.A2672.109119}

\bibitem{brewer2015tough}
Robert S. Brewer, G. E. Xu, T. K. Lee, M. Katchuck, C. A. Moore, and P. M. Johnson. 2015. Tough Shift: Exploring the Complexities of Shifting Residential Electricity Use Through a Casual Mobile Game. In \emph{Proceedings of the 2015 Annual Symposium on Computer-Human Interaction in Play (CHI PLAY '15)}. ACM, 307--317. DOI: \href{https://doi.org/10.1145/2793107.2793108}{https://doi.org/10.1145/2793107.2793108}

\bibitem{kjeldskov2015ecoforecasting}
Jesper Kjeldskov, Mikael B. Skov, Jeni Paay, Dennis Lund, Tue Madsen, and Michael Nielsen. 2015. Eco-Forecasting for Domestic Electricity Use. In \emph{Proceedings of the 33rd Annual ACM Conference on Human Factors in Computing Systems (CHI '15)}. ACM, 1985--1988. DOI: \href{https://doi.org/10.1145/2702123.2702318}{https://doi.org/10.1145/2702123.2702318}

\bibitem{koide2019carbon}
Ryota Koide, Michael Lettenmeier, Satoshi Kojima, Viivi Toivio, Aryanie Amellina, and Lewis Akenji. 2019. Carbon footprints and consumer lifestyles: an analysis of lifestyle factors and gap analysis by consumer segment in Japan. \emph{Sustainability} 11, 21 (2019), 5983. DOI: \href{https://doi.org/10.3390/su11215983}{https://doi.org/10.3390/su11215983}

\end{thebibliography}

\end{document}